\section{Smearing estimates and conditional moments}\label{sec:moments}

The spectral and pin-motion arguments require bounds for smooth averages,
including averages conditioned on events whose probabilities may vanish
with the mesh.  We first obtain these bounds in flat and bounded-boundary
domains, then in the plane, and finally under separated pin obstacles.
The same midpoint amplification is used at each stage, with the offset
parities checked explicitly in the last application.

A lattice smearing of radius $a\geq1$ is a statistic
$Y=\sum_x w_xH(x)$ whose weights are supported in a ball of radius $a$
and satisfy
\begin{equation}\label{eq:normalized-weights}
 \sum_x\abs{w_x}\leq A,\qquad
 \sup_x\abs{w_x}\leq A a^{-2}.
\end{equation}
The constant $A$ is fixed in all uniform assertions.  Signed weights are
allowed.  A \emph{zero-total} smearing also has $\sum_xw_x=0$, so it is
independent of the choice of height anchor.  Affine interpolation of a
bounded test density of size $O(a^{-2})$ gives weights of this form, with
the same total mass as the test.  At mesh $\delta$, a fixed physical
radius is divided by $\delta$ in these lattice-unit definitions.

\subsection{An amplification lemma for uniform arrays}

The rounded midpoint inequality originates in
Klartag--Lehec \cite[Theorem~1.4]{KlartagLehec2019}.  We use its
nonnegative-function formulation in
Gozlan--Roberto--Samson--Tetali \cite[Theorem~3]{GRST2019}.
In one dimension it states that, if finitely supported nonnegative
functions on $\Z$ satisfy
\[
 f(x)g(y)\leq
 u\bigl(\lfloor(x+y)/2\rfloor\bigr)
 v\bigl(\lceil(x+y)/2\rceil\bigr),
\]
for all $x,y\in\Z$, then $(\sum f)(\sum g)\leq(\sum u)(\sum v)$.

The nonnegative formulation, including functions with zeros, is essential
because our applications use indicators.  Here is a direct justification
in the finite-support case from the entropy inequality in
\cite[Theorem~8, equation~(9)]{GRST2019}.  Write $F_0=\sum f$,
$F_1=\sum g$, $U_-=\sum u$, and $U_+=\sum v$.  If either $F_i$ is zero
the conclusion is immediate.  Otherwise normalize $f,g$ to probability
laws $\nu_0,\nu_1$ and use their monotone coupling.  Let $\nu_-,\nu_+$
be the laws of its rounded-down and rounded-up midpoints, and put
$\mathcal H(\nu)=\sum_z\nu(z)\log\nu(z)$.  The cited entropy inequality
says
\[
 \mathcal H(\nu_-)+\mathcal H(\nu_+)
 \le \mathcal H(\nu_0)+\mathcal H(\nu_1).
\]
Every pair with positive coupling mass has positive input values, so the
pointwise midpoint hypothesis forces both corresponding output values to
be positive.  Logarithms may therefore be integrated on these finite
supports, without modifying any zero value.  The pointwise hypothesis and
the finite log-sum inequality give
\begin{align*}
 \mathcal H(\nu_0)+\mathcal H(\nu_1)+\log(F_0F_1)
 &\le \sum_z\nu_-(z)\log u(z)+\sum_z\nu_+(z)\log v(z)\\
 &\le \mathcal H(\nu_-)+\mathcal H(\nu_+)+\log(U_-U_+).
\end{align*}
Combining these two inequalities proves the scalar assertion.  The same
inequality holds for arrays on $\Z^d$, with coordinatewise rounding.
Indeed, for each fixed pair of last coordinates,
apply the lower-dimensional assertion to the slices; then apply the scalar
assertion to their sums.  This also explains why output supports need not
be the same as the input supports.

The next lemma turns a central probability bound into moments.  It allows
the parity classes needed for pinned heights, while requiring midpoint
admissibility only within each class.

\begin{lemma}[Midpoint tail amplification]\label{lem:midpoint-tails}
Let $\mathcal S$ be a finite set of integer arrays with its uniform law,
let $Y$ be a real linear statistic, and let $M>0$.  Partition $\mathcal S$ into at
most $m_*$ classes.  Suppose that whenever $x,y$ belong to the same class,
both coordinatewise rounded midpoints belong to $\mathcal S$, and that
their $Y$ values differ from $(Y(x)+Y(y))/2$ by at most $M$.  Suppose also
that, for every class in use,
\begin{equation}\label{eq:class-central-mass}
 \PP\{\abs Y\leq M,\ x\hbox{ belongs to that class}\}\geq c>0.
\end{equation}
Then, for constants depending only on $c$ and $m_*$,
\[
 \PP(\abs Y>tM)\leq C e^{-c't}\quad(t\geq1),\qquad
 \norm{Y}_{L^p}\leq C_pM\quad(p<\infty).
\]
The rounded outputs need not belong to the input class.
\end{lemma}

\begin{proof}
Consider the enlarged positive dyadic intervals
$I_k=[2^kM-4M,2^{k+1}M+4M]$, and put $p_k=\PP(Y\in I_k)$.
Choose a class carrying at least $p_k/m_*$ of this mass.  Pair its
indicator on $I_k$ with its indicator on $[-M,M]$.  Both rounded output
statistics belong to $I_{k-1}$: the half-width error is at most
$2M+M/2+M<4M$.  Apply the product midpoint inequality with both output
arrays equal to the indicator of $\mathcal S\cap\{Y\in I_{k-1}\}$.
After division by $\abs{\mathcal S}^2$, it gives
\begin{equation}\label{eq:dyadic-tail-recursion}
 p_{k-1}^2\geq c p_k/m_*.
\end{equation}
The admissibility assumption is needed only where the two input
indicators are nonzero; no output class restriction occurs.

For $k\geq4$ these enlarged intervals have bounded overlap.  Hence, in
an initial range of bounded length depending only on $c,m_*$, some $p_{k_0}$
is at most $c/(2m_*)$.  Write $b=m_*/c$.  Iterating
\eqref{eq:dyadic-tail-recursion} gives
$b p_{k_0+j}\leq2^{-2^j}$.  The bounded size of the initial range absorbs
the earlier terms, so $p_k\leq C\exp(-c'2^k)$ for all $k$.  Repeat with
negative intervals and sum the resulting tail bounds.  Integrating the
tail proves the moment assertion.
\end{proof}

Integer Lipschitz arrays are stable under these two midpoint operations:
an integral edge bound on the averaged coordinates is preserved by taking
floors at both endpoints or ceilings at both endpoints.  A fixed integral
boundary or anchor is preserved as well.  For a statistic with weights
$w_x$, the rounding error is at most $\frac12\sum_x\abs{w_x}$.  These
observations verify the lemma's admissibility hypotheses in the unpinned
applications.

\subsection{Bounded-boundary height problems}

\begin{lemma}[Loop counts in a distance shell]\label{lem:loop-shells}
In a finite flat-boundary height problem, fix an interior site $x$.
The number of level loops surrounding $x$ whose graph inradii lie between
$d$ and $2d$ has moments bounded at every order, uniformly in the domain
and in $d\geq1$.  The same holds when the shell is defined by Euclidean
closest approach to $x$.  Both assertions also hold in a bounded face of
a connected constant-height skeleton.
\end{lemma}

\begin{proof}
Choose $\eta>0$ small enough that $\rho(1+\eta)<1$, with slack for bounded
lattice rounding at large distances.  Divide $[d,2d]$ into a bounded
number of intervals with ratio at most $1+\eta$.  Expose surrounding loops
from the exterior, and consider the first one whose graph inradius
lies in one such interval $[d', (1+\eta)d']$.  If it exists, its interior
has the empty-exterior loop law.  If its graph inradius is $n$, then
$n\le(1+\eta)d'$.  Every further loop with graph inradius at least $d'$
surrounds $\Lambda_{\lfloor d'\rfloor}(x)$.  For sufficiently large $d'$,
\[
 \lceil\rho n\rceil<\lfloor d'\rfloor,
\]
because $\rho(1+\eta)<1$.  Thus two such loops imply the exact event
in the imported upper estimate, in the same graph convention.  Their
conditional probability is at most $1-c_0$.
Reapply this argument after each successive pair of exposed loops.
The count in the interval is dominated by one plus twice a geometric
random variable with a fixed success parameter.  A bounded number of
such bounds handles $[d,2d]$.

At bounded distances, disjoint loops approaching that close must use
distinct vertices in a fixed finite ball, so their number is bounded
deterministically.  This handles the rounding exceptions.  No imported
upper bound was applied to the original polygon or skeleton face; every
application followed exposure of a surrounding loop.  The flat loop
description in \cref{lem:gray} therefore suffices for the stated domains.

Finally, triangular graph distance and Euclidean distance are comparable
by fixed constants, and a contour is within bounded distance of its
incident face centers.  A Euclidean closest-approach shell $[d,2d]$ is
therefore covered by a fixed finite number of graph-inradius shells,
with bounded small-radius exceptions.  Summing the graph-shell moment
bounds proves the Euclidean assertion.  This metric comparison is made
only after the strict graph-radius contraction has been proved.
\end{proof}

\begin{proposition}[Interior smearing moments]\label{prop:smearing-moments}
Consider a uniform height problem with boundary values bounded in absolute
value by $K$.  Let the center of a radius-$a$ smearing have distance $n$
from the boundary, with $1\leq a\leq c_0n$ for a sufficiently small fixed
$c_0>0$.  Then, for every finite $p$,
\begin{equation}\label{eq:smearing-moments}
 \norm{\textstyle\sum_xw_xH(x)}_{L^p}
 \leq C_p\bigl(1+\sqrt{\log(n/a)}\bigr).
\end{equation}
Here $C_p$ depends on $p,A,K$, but not on boundary regularity.  The same
statement holds for bounded data on the inner face chain of a dual cut,
and in the bounded skeleton faces considered above.
\end{proposition}

\begin{proof}
First put zero on the boundary.  Conditional on all loop geometry,
independent centered jump signs give
\begin{equation}\label{eq:common-loop-covariance}
 \EE[H(x)H(y)]
 =\EE\bigl[\#\{\hbox{loops surrounding both }x\hbox{ and }y\}\bigr].
\end{equation}
Let $s=\max(1,\dist(x,y))$, for $x,y$ in the smearing ball.  Their boundary
distances are comparable to $n$.  A loop surrounding both points and
approaching $x$ within distance $2d\ll s$ supplies a wall arm from scale
$O(d)$ to scale of order $s$.  The contour changes only one spin color
along its entire length.  A surrounding double circuit of that color
blocks the arm: at a vertex of the double circuit all three adjacent
spins have that color's fixed value, so its wall cannot cross there.

Flat height boundary conditions fix both colors on a connected boundary
chain.  For either color used as the first spin, constancy of the other
color on this chain is an $E_0$ requirement; its single label does not
change the first marginal.  The buffers between scales $d$ and $s$ lie
away from the boundary.  Thus \cref{lem:buffer-circuits}, for either
possible wall color, bounds the arm probability by $C(d/s)^\alpha$.
Together with the second moment in \cref{lem:loop-shells}, Cauchy--Schwarz
bounds the expected common-loop count in this shell by
$C(d/s)^{\alpha/2}$.  Summing the small shells is finite.  Each shell from
scale $s$ up to scale $O(n)$ contributes at most a constant.  There are no
larger enclosing shells, since an enclosing loop cannot surround a fixed
boundary site.  Consequently
\begin{equation}\label{eq:point-covariance-upper}
 0\leq\EE[H(x)H(y)]\leq C\bigl(1+\log(n/s)\bigr).
\end{equation}

Summing \eqref{eq:point-covariance-upper} against $\abs{w_xw_y}$ yields
\begin{equation}\label{eq:smearing-variance}
 \EE\Bigl(\sum_xw_xH(x)\Bigr)^2
 \leq C\bigl(1+\log(n/a)\bigr).
\end{equation}
Indeed the part $\log(n/a)$ is controlled by the total weight.  For the
remaining singularity, graph balls contain $O(r^2)$ sites, so the bounds
in \eqref{eq:normalized-weights} give a uniform average of
$\log(a/\max(1,\dist(x,y)))$.  The one-site contribution is bounded by
$Ca^{-2}\log a$.  This calculation also covers signed weights.

Chebyshev's inequality gives a fixed central probability on the scale
$M=C(1+\sqrt{\log(n/a)})$, chosen also to dominate the rounding error.
The height arrays with zero boundary have admissible rounded midpoints,
so \cref{lem:midpoint-tails} with one class proves
\eqref{eq:smearing-moments} in the flat case.

For bounded boundary data, compare with the zero-boundary height shifted
by $-\lceil K\rceil$ and $+\lceil K\rceil$.  A nonnegative smearing is
stochastically trapped between the corresponding shifted smearings.
Its two tails are therefore bounded by those of the zero-boundary problem,
with an additional deterministic error at most $A\lceil K\rceil$.
Split signed weights into positive and negative parts and use the triangle
inequality in $L^p$.  This proves the stated bound.  The proof uses only
the flat loop representation and the incident-star convention, and hence
also applies to the asserted dual-cut and skeleton boundaries.
\end{proof}

In particular, a normalized average on a ball whose radius is a fixed
small fraction of its distance to the boundary has moments bounded
uniformly as that distance tends to zero, with mesh sent to zero first.
This observation will be used at both ordinary boundary points and the
two marks.

\begin{proposition}[Plane zero-total moments]\label{prop:plane-moments}
For every zero-total smearing satisfying \eqref{eq:normalized-weights},
\begin{equation}\label{eq:plane-zero-total-moments}
 \norm{\textstyle\sum_x w_xH(x)}_{L^p(\mu_\infty)}\leq C_p,
 \qquad p<\infty,
\end{equation}
uniformly in its radius, location, and lattice rotation.  The same bound
holds in sufficiently larger finite free volumes, for a height anchored
outside the construction.
\end{proposition}

\begin{proof}
Search from the exterior for a double-$B$ circuit in an annulus between
$Ka$ and $2Ka$ about the support, where $K$ is a fixed sufficiently large
constant.  \Cref{lem:buffer-circuits} gives probability at least $c>0$,
uniformly in $a$ and in a still larger free volume.  The search does not
inspect the interior.  Its inner height trace lies in one pair of
consecutive integers by \cref{lem:double-cut}.  Subtract the lower integer
of that pair.  Conditional on the trace, the interior is a bounded-boundary
height problem.  Its boundary distance from the support is comparable to
$a$, and $K$ may be chosen to meet the small-fraction requirement in
\cref{prop:smearing-moments}.  That proposition bounds the second moment
of the shifted smearing by a constant.  The shift does not alter the
original smearing, since its total is zero.  We have therefore obtained
an unconditional central probability $\PP(\abs Y\leq M)\geq c'>0$ for a
constant $M$.

The free pair increment law is the uniform anchored-height law, as noted
after \cref{prop:plane-mixing}.  Its arrays have admissible rounded
midpoints and bounded rounding error.  Apply \cref{lem:midpoint-tails}
in this finite volume, with the anchor outside the searched circuit.
This gives the bound for every $p$, uniformly in that volume and in $a$.
For each fixed lattice smearing, let the volume expand and use
\cref{prop:plane-mixing}; the bounds pass to the limit.  The constants are
uniform in $a$ and the listed symmetries, proving the assertion.
\end{proof}

\subsection{An isolated component switch}

The following finite identity has two uses: it supplies the parity classes
needed for conditional moments, and it prevents the limiting covariance
coefficient from vanishing.

\begin{lemma}[A fair switch between opposite cuts]\label{lem:component-switch}
Suppose that two nested double-$B$ circuits have opposite signs, and that
their intervening collar contains no $E_0$ edge or pinned $W$ site.  There
is a component of $B$-changing edges confined to this collar whose fair
$W$ sign can be flipped so as to change every height inside the inner
cut by $4$ in absolute value relative to every height outside the outer
cut.  The flip changes no spin pin outside the collar.  It preserves any inner
height prescription invariant under a common additive shift in $4\Z$,
while leaving exterior heights fixed.

Conditional on $B$ and the other $W$ component labels, the two choices
are equally likely.  If a zero-total statistic has total weight $\lambda$
inside the inner cut, no weight in the collar, and all its other weight
outside the outer cut, its two values differ by $4\abs\lambda$ and its
conditional variance is at least $4\lambda^2$.
\end{lemma}

\begin{proof}
On every oriented edge the Gray increment gives the identity
\begin{equation}\label{eq:component-potential-identity}
 d(H+BW/2)=W\,dB.
\end{equation}
Let $J$ be one connected component of $B$-changing edges, and put
$\omega_J=\one_JdB$, extended by zero to other edges.  A triangle has
either no $B$-changing edge or two incident ones in the same component.
Thus $\omega_J$ has zero curl and has a potential $\psi_J$ on the simply
connected ambient graph.  On the vertices of $J$ this potential agrees
with $B$ up to an additive constant.  Every other vertex can be connected
to $J$ by a path with no edge of $J$ before its endpoint, and so its
potential continues one of the two values on $J$.  The potential has
exactly two values, differing by two.

No component of $B$-changing edges can cross a double-$B$ cut.  Between
the two cuts, the sum of the potential changes of all such components,
from the outer cut to the inner one, is the total change of $B$, namely
$2$ or $-2$.  At least one component $J$ in the collar therefore has a
nonzero change.  Its potential is constant on the inner region and on
the outer region, with a difference of two.  Normalize it to zero outside
the outer cut.

The collar is free of $E_0$ and pinned $W$ sites.  Consequently the
$W$-equality component containing $J$ consists precisely of its vertices,
and its common value $w$ is an independent fair sign by
\cref{lem:spin-fkg}.  Flip this sign and leave all other $W$ components
unchanged.  Formula \eqref{eq:component-potential-identity} gives, at every
vertex not in $J$,
\[
 H'-H=-2w\psi_J,
\]
when the height is anchored outside the outer cut.  The correction from
$BW/2$ vanishes at these vertices.  Thus the height change is zero
outside and is $+4$ or $-4$ throughout the inner region.  The spin pins
outside the collar are untouched.  The two statistic values differ by
$4\abs\lambda$, so a fair choice between them has variance
$(4\lambda)^2/4=4\lambda^2$.
\end{proof}

For two fixed disjoint test disks, \cref{lem:buffer-circuits} can require
these opposite cuts around one disk while avoiding the other, with a
uniform positive probability.  Applied to the difference of their
unit-mass averages, \cref{lem:component-switch} therefore gives a strictly
positive, scale-uniform second-moment lower bound.  This assertion is
also uniform in a large finite-volume approximation to the plane.

\subsection{Moments conditional on pin obstacles}

For the following moment statement, impose one extra condition on the
separated obstacles of \cref{prop:pin-product}.  We call a pin pattern
\emph{interval-compatible} if, along a connected chain on obstacle $i$,
its exact height description is
\begin{equation}\label{eq:interval-pin-description}
 H(v)\in I_{i,v}+4t_i\quad(v\hbox{ on the chain}),\qquad t_i\in\Z,
\end{equation}
where the $I_{i,v}$ are prescribed integer intervals and at least one is
a singleton.  The spin event must be exactly this requirement for some
$t_i$, with no additional height restriction.  The patterns used later
have intervals of one or two consecutive integers.  They are represented
by $B$ pins and by $W=+1$ on a bounded number of connected subpaths,
as in \cref{prop:pin-product}.  A singleton makes each $t_i$ unique after
a common lift has been chosen.  Reflections and translations preserve
this definition.

\begin{proposition}[Conditional obstacle moments]\label{prop:obstacle-moments}
Fix finitely many separated interval-compatible obstacles with filled
isolating neighborhoods $U_i$ as in \cref{prop:pin-product}.  Let a test
disk lie in a pin-free simply connected neighborhood $U_0$, with room
for a surrounding buffered cut whose filled interior lies in $U_0$.
Require $\overline U_0$ to be disjoint from every $\overline U_i$.
Let $Y$ be a zero-total smearing in the test disk, satisfying
\eqref{eq:normalized-weights} at its scale.  Conditional on all the pins,
\begin{equation}\label{eq:obstacle-moment-bound}
 \norm Y_{L^p}\leq C_p\qquad(p<\infty).
\end{equation}
The constants depend on the macroscopic arrangement, the weight bound,
the number of obstacles, and the bound on W-subpath counts, but not on
the mesh or the microscopic pin probabilities.  No uniformity as obstacle
separations vanish is asserted here.
\end{proposition}

\begin{proof}
If the obstacle family is empty, use \cref{prop:plane-moments} directly.
For a nonempty family, begin in a finite free volume and anchor the
height at a singleton on obstacle 1 with its prescribed unshifted
value, so $t_1=0$.  The conditional
pair law becomes uniform counting of admissible anchored height arrays
satisfying \eqref{eq:interval-pin-description}.  Partition them by the
parities of $t_2,\ldots,t_m$.  In a pair of arrays in the same class,
$(t_i+t_i')/2$ is an integer.  Both rounded midpoint arrays therefore
satisfy the same interval requirements with this new offset: the interval
endpoints are integral.  The edge bounds and the anchor are preserved.
The rounding error of $Y$ is bounded by $A/2$.  We verify the remaining
central-mass hypothesis of \cref{lem:midpoint-tails}.

First retain W-constancy as $E_0$ and temporarily drop its plus labels.
Around the test find a double-$B$ cut by an exterior stopping search.
Around each non-anchor obstacle find two nested double-$B$ cuts of
opposite signs in its pin-free collar.  Choose all these search buffers
disjoint.  \Cref{lem:buffer-circuits} gives their joint event $\mathcal A$
probability bounded below, conditional on the $B$ pins and $E_0$.
Restoring the bounded number of W-plus labels loses at most the fixed
factor explained after \cref{lem:double-cut}.  Thus
\begin{equation}\label{eq:obstacle-cut-probability}
 \PP(\mathcal A\mid\hbox{all pins})\geq c_1>0.
\end{equation}

The test search has not revealed its strict interior.  Conditional on
its trace, its height boundary lies in one consecutive pair, and there
are no interior pin constraints.  Subtract that pair's lower value.
The zero total removes the shift, and \cref{prop:smearing-moments} supplies
a uniform conditional second moment.  If necessary cover the test support
by a fixed number of smaller balls to meet that proposition's small-radius
condition.  For a sufficiently large fixed $M$, Chebyshev gives
\begin{equation}\label{eq:obstacle-central-event}
 \PP(\mathcal A,\abs Y\leq M\mid\hbox{all pins})\geq c_2>0.
\end{equation}
Here we conditioned on the searches and their exterior traces, not on
all the $B$ values inside the test disk.

Now expose all $B$ values and use \cref{lem:component-switch} separately
in each non-anchor obstacle collar.  Choose a contributing component
measurably from $B$.  Its $W$ label is an independent fair bit, unaffected
by the required plus labels because the collar has no pinned W sites or
$E_0$ edges.  Flipping it changes $t_i$ by an odd integer and changes no
other offset.  The anchor, every other obstacle, and the whole test neighborhood lie
outside $U_i$, hence outside this collar's outer cut.  The flip consequently
preserves $Y$, all other offsets, all pins, and $\mathcal A$.
Distinct collars give distinct independent bits.

Conditional on $B$ and all remaining W-component labels, these switches
act transitively on the $2^{m-1}$ parity vectors, preserve probability,
and preserve the event in \eqref{eq:obstacle-central-event}.  Each parity
class thus receives at least $2^{-(m-1)}c_2$ of that central probability.
This is \eqref{eq:class-central-mass}, with a number of classes independent
of the mesh.  \Cref{lem:midpoint-tails} proves every moment bound.

Finally let the free volume expand at each fixed mesh.  The finite pin
event has positive plane probability by \cref{prop:plane-mixing}; cylinder
convergence therefore passes to its conditional law.  The uniform moment
bounds pass as well, proving \eqref{eq:obstacle-moment-bound} in the plane.
\end{proof}

\begin{corollary}[Equal-offset moments with varying separations]
\label{cor:equal-offset-moments}
For a nonempty obstacle family in the preceding setup, impose the additional event
$\mathcal E=\{t_i=t_1\hbox{ for every }i\}$.  The obstacles may now vary
or approach one another, while staying away from a fixed test disk.
Fix an annular buffer around that disk, with positive clearance from
the disk and with fixed inner and outer radii.  Suppose the buffer and
its filled interior contain no pins.
Suppose that, under conditioning on the pins and $\mathcal E$, an exterior
search finds a surrounding double-$B$ cut in this buffer, with probability
at least $c>0$ uniformly in the varying arrangements.  Suppose that
$\mathcal E$ can be checked in the exterior of the selected cut, so its
strict interior remains free given the trace.  Then the zero-total test
in that disk has moments bounded uniformly in those arrangements.

A sufficient condition is that the cut probability conditional on the
pins alone is at least $c_0>0$, that
$\PP(\mathcal E\mid\hbox{pins})\to1$ uniformly along the varying family,
and that the same exterior measurability holds.
\end{corollary}

\begin{proof}
Anchor at obstacle 1.  Equal offsets make all the offsets zero, so the
admissible height arrays now satisfy fixed integer interval constraints.
Their rounded midpoints are admissible without a parity partition.
On the asserted cut event, the trace lies in a consecutive pair and the
conditional interior is free of further constraints.  The fixed buffer
bounds its distance from the test above and below, uniformly in the
arrangement.  Covering the support by a fixed number of smaller balls
if necessary, \cref{prop:smearing-moments} therefore gives the uniform
central probability in \eqref{eq:obstacle-central-event}.  Apply
\cref{lem:midpoint-tails} with one class.  This proof uses
neither the separate obstacle collars nor bounds on moments outside the
equal-offset event.

For the sufficient condition, if
$\eps=\PP(\mathcal E^c\mid\hbox{pins})$, then the conditional cut probability
after adding $\mathcal E$ is at least $(c_0-\eps)/(1-\eps)$.
This is bounded below for all sufficiently small $\eps$.  Exterior
measurability ensures that the conditioning has not changed the free
interior once the trace is fixed.  Finite-volume approximation justifies
the counting argument exactly as in \cref{prop:obstacle-moments}.
\end{proof}
